\documentclass{article}

\usepackage{amsmath}
\usepackage{amssymb}
\usepackage{polski}

\usepackage[utf8]{inputenc}
\usepackage[T1]{fontenc}

\usepackage{listings}
\lstset{
    language=SQL,
    inputencoding=utf8x,
    extendedchars=\true,
    literate={ą}{{\k{a}}}1
             {Ą}{{\k{A}}}1
             {ę}{{\k{e}}}1
             {Ę}{{\k{E}}}1
             {ó}{{\'o}}1
             {Ó}{{\'O}}1
             {ś}{{\'s}}1
             {Ś}{{\'S}}1
             {ł}{{\l{}}}1
             {Ł}{{\L{}}}1
             {ż}{{\.z}}1
             {Ż}{{\.Z}}1
             {ź}{{\'z}}1
             {Ź}{{\'Z}}1
             {ć}{{\'c}}1
             {Ć}{{\'C}}1
             {ń}{{\'n}}1
             {Ń}{{\'N}}1
}

\title{Liczby Zespolone - Zestaw I}
\date{2026}
% \author{Maciej Różewicz}

\begin{document}
\maketitle

\begin{enumerate}
    \item Udowodnić następujące własności działań w~zbiorze liczb zespolonych:
    \begin{enumerate}
        \item $ \forall z_{1}, z_{2} \in \mathbb{C}: z_{1} + z_{2} = z_{2} + z_{1} $,
        \item $ \forall z_{1}, z_{2}, z_{3} \in \mathbb{C}: \left( z_{1} + z_{2} \right) +z_{3} = z_{1} + \left( z_{2} + z_{3} \right) $,
        \item $ \forall z \in \mathbb{C}: z + \textbf{0} = z, gdzie\, \textbf{0} = (0, 0) $,
        \item $ \forall z_{1}, z_{2} \in \mathbb{C}: z_{1}z_{2} = z_{2}z_{1} $,
        \item $ \forall z_{1}, z_{2}, z_{3} \in \mathbb{C}: \left(z_{1}z_{2}\right)z_{3} = z_{1}\left(z_{2}z_{3}\right) $,
        \item $ \forall z \in \mathbb{C}-\{\textbf{0}\} \, \exists! \tilde{z} \in \mathbb{C}: z\tilde{z}=\textbf{1}  $%-\left{\textbf{0}\right} \exists! \tilde{z} \in \mathbb{C}: z\tilde{z}=\textbf{0}, gdzie \textbf{1} = (1, 0) $
        \item $ \forall z_{1}, z_{2}, z_{3} \in \mathbb{C}: z_{1} \left(z_{2} + z_{3}\right) = z_{1}z_{2} + z_{1}z_{3}$,
    \end{enumerate}
    
    \item Pokazać, że zbiór liczb zespolonych $\mathbb{C}$ z~działaniami zdefiniowanymi w~następujacy sposób:
    \begin{itemize}
        \item $ \left( a_{1}, a_{2} \right) + \left( b_{1}, b_{2} \right) = \left( a_{1} + b_{1}, a_{2} + b_{2} \right) $
        \item $ \left( a_{1}, a_{2} \right) \bullet \left( b_{1}, b_{2} \right) = \left( a_{1}b_{1} - a_{2}b_{2}, a_{1}b_{2} + a_{2}b_{1} \right) $
    \end{itemize}
    jest ciałem przemiennym.
    
    \item Niech $z_{1}, z_{2}, z \in \mathbb{C}$. Podać interpretację geometryczną dla:
    \begin{enumerate}
        \item $z$,
        \item $\bar{z}$,
        \item $z_{1}+z_{2}$,
        \item $z_{1}-z_{2}$,
        \item $|z|$,
        \item * $z_{1}z_{2}$
    \end{enumerate}

    \item Niech $z_{1}, z_{2}, z \in \mathbb{C}$. Uzasadnić zależności:
    \begin{enumerate}
        \item $\bar{z_{1}z_{2}} = \bar{z_{1}}\bar{z_{2}}$,
        \item $\bar{\left( \frac{z_{1}}{z_{2}} \right)}  = \frac{\bar{z_{1}}}{\bar{z_{2}}}, z_{2} \neq \textbf{0}$,
        \item $|z_{1}z_{2}| = |z_{1}||z_{2}|$,
        \item $|\frac{z_{1}}{z_{2}}| = \frac{|z_{1}|}{|z_{2}|}, z_{2} \neq \textbf{0}$,
        \item $\bar{z_{1} + z_{2}} = \bar{z_{1}} + \bar{z_{2}}$,
        \item $\bar{z_{1} - z_{2}} = \bar{z_{1}} - \bar{z_{2}}$,
        \item $z + \bar{z} = 2Re(z)$,
        \item $z - \bar{z} = 2i Im(z)$,
    \end{enumerate}

    \item Wykonać obliczenia:
    \begin{enumerate}
        \item $(3-i) + (2+4i)$,
        \item $(2-3i) - (5i)$,
        \item $(1-i)(2+4i)$,
        \item $(2+3i)(2-3i)$,
        \item $\frac{2-3i}{3+i}$,
        \item $\frac{4+2i}{1-2i}$,
        \item $\frac{3+2i}{i}$,
        \item $\frac{1-2i}{3i+2}$,
        \item $\frac{1}{i} + \frac{2+i}{2-i}$,
        \item $\frac{(1+i)^n}{(1-i)^{n-2}}, n \in \mathbb{N}$,
        \item $(\cos{\alpha} + i\sin{\alpha})^n, n \in \mathbb{N}$
    \end{enumerate}

    \item Rozwiązać równanie, dla $z\in\mathbb{C}$ i $x, y\in\mathbb{R}$:
    \begin{enumerate}
        \item $(1 + 2i)x + (2-2i)y = 5 + 4i$,
        \item $z^{2} + (3-2i)z + 1 -3i = 0$,
        \item $(z-2)^{2} + 2 Im(z) + z = 1$,
        \item $|z| + z = 1 + i$,
        \item $\frac{1+i}{z} = \frac{2-3i}{\bar{z}}$,
        \item $2z + (3-i)\bar{z} = 4 + 4i$,
        \item $z + i = \bar{z+i}$,
        \item $z\bar{z} + (z - \bar{z}) = 3 + 2i$
    \end{enumerate}

    \item Przedstawić graficzną reprezentację dla równania lub nierówności dla $z, z_{j} \in \mathbb{C}$:
    \begin{enumerate}
        \item $|z-z_{0}| = r$,
        \item $|z-z_{0}| \leq r$,
        \item $|z-z_{0}| < r$,
        \item $|z-z_{0}| \geq r$,
        \item $|z-z_{0}| > r$,
        \item $r \leq |z-z_{0}| \leq R$,
        \item $|z-z_{1}| = |z-z_{2}|$,
        \item $|z-z_{1}| < |z-z_{2}|$,
        \item $|z-z_{1}| > |z-z_{2}|$
    \end{enumerate}

    \item Zaznaczyć na płaszczyźnie zespolonej liczby zespolone $z$, dla których:
    \begin{enumerate}
        \item $\frac{z+4}{z-2i}$ jest liczbą rzeczywistą,
        \item $\frac{z}{iz + 4}$ jest czysto urojona,
        \item $|z+i|=3$,
        \item $|\frac{z-3}{z-3i}|>1$,
        \item $2 < |z+2-i|\leq 3$,
        \item $|\bar{z} + 2 -i| \leq |z|$
    \end{enumerate}

    \item Podane liczby zespolone przedstawić w~postaci trygonometrycznej:
    \begin{enumerate}
        \item $\sqrt{3} + i$,
        \item $1 + \sqrt{3}i$,
        \item $-\sqrt{3} + i$,
        \item $-1$,
        \item $1 + i$,
        \item $-\frac{1}{2} - \frac{\sqrt{3}}{2}i$,
        \item $3 + i$,
        \item $4 - i$,
        \item $-2 + i$
    \end{enumerate}

    \item Wykorzystując postać trygonometryczną liczby zespolonej wykonać działania:
    \begin{enumerate}
        \item $(1-i)(\sqrt{3}+i)$,
        \item $(4 + 4i)(-3+3i)$,
        \item $\frac{2+2i}{1-i}$,
        \item $\frac{3i}{1+i}$
    \end{enumerate}

\end{enumerate}

\newpage
\textbf{Odpowiedzi:}
\begin{description}
    \item [5:] a)$5+3i$, b)$2-8i$, c)$6+2i$, d)$3$, e)$\frac{3}{10}-\frac{11}{10}i$, f)$2i$, g)$2-3i$, h)$\frac{4}{13}-\frac{7}{13}i$, i)$\frac{3}{5}+\frac{9}{5}i$, j)$2i^{n-1}$, k)$\cos{(n\alpha)} + \sin{(n\alpha)}i$
    \item [6:] a)$3+i$, b)$-2+i$, $2+i$, c)$\frac{3}{2} + (1+\frac{\sqrt{7}}{2}), \frac{3}{2} + (1-\frac{\sqrt{7}}{2})i$, d)$i$, e)$\empty$, f)$-4i$, g)$x-i, x\in\mathbb{R}$, h)$\sqrt{2}+i, -\sqrt{2}+i$
    \item [7:] a)$ \{ (x,y): (x-x_{0})^2 + (y-y_{0})^2 = r^2 \} $, b)$ \{ (x,y): (x-x_{0})^2 + (y-y_{0})^2 \leq r^2 \} $, c)$ \{ (x,y): (x-x_{0})^2 + (y-y_{0})^2 < r^2 \} $, d)$ \{ (x,y): (x-x_{0})^2 + (y-y_{0})^2 \geq r^2 \} $, e)$ \{ (x,y): (x-x_{0})^2 + (y-y_{0})^2 > r^2 \} $, f)$ \{ (x,y): (x-x_{0})^2 + (y-y_{0})^2 \geq r^2 \land (x-x_{0})^2 + (y-y_{0})^2 \leq R^2  \} $, g)$ \{ (x,y): y = -x\frac{x_{2}-x_{1}}{y_{2}-y_{1}} + \frac{x_{2}^2-x_{1}^2+y_{2}^{2}-y_{1}^2}{2(y_{2}-y_{1})} \} $, h)$ \{ (x,y): y < -x\frac{x_{2}-x_{1}}{y_{2}-y_{1}} + \frac{x_{2}^2-x_{1}^2+y_{2}^{2}-y_{1}^2}{2(y_{2}-y_{1})} \} $, i)$ \{ (x,y): y > -x\frac{x_{2}-x_{1}}{y_{2}-y_{1}} + \frac{x_{2}^2-x_{1}^2+y_{2}^{2}-y_{1}^2}{2(y_{2}-y_{1})} \} $, 
    \item [8:] a)$\{ (x,y): y = \frac{1}{2}x + 2,\,x\neq0 \}$, b)$\{ (x, y): x=0\,y\neq4 \}$, c)$\{ (x, y): x^{2}+(y+1)^{2} = 9 \}$, d)$ \{ (x, y): y>x \}$, e)$\{ (x,y): (x+2)^{2} + (y-1)^{2} > 4 \land (x+2)^{2} + (y-1)^{2} \leq 9 \}$, f)$ \{ (x,y): y < -2x - \frac{3}{2} \}$
    \item [9:] a)$2(\cos{\frac{\pi}{6}} + i\sin{\frac{\pi}{6}})$, b)$2(\cos{\frac{\pi}{3}} + i\sin{\frac{\pi}{3}})$, c)$2(\cos{\frac{5\pi}{6}} + i\sin{\frac{5\pi}{6}})$, d)$1(\cos{\pi} + i\sin{\pi})$, e)$\sqrt{2}(\cos{\frac{\pi}{4}} + i\sin{\frac{\pi}{4}})$, f)$1(\cos{\frac{4\pi}{3}} + i\sin{\frac{4\pi}{3}})$, g)$\sqrt{10}(\cos{(\arctan{\frac{1}{3}})} + i\sin{(\arctan{\frac{1}{3}})})$, b)$\sqrt{17}(\cos{(\arctan{\frac{-1}{4}})} + i\sin{(\arctan{\frac{-1}{4}})})$, b)$\sqrt{5}(\cos{(\pi - \arctan{\frac{1}{2}})} + i\sin{(\pi - \arctan{\frac{1}{2}})})$
    \item [10:] a)$2\sqrt{2}(\cos{(-\frac{\pi}{12})} + i\sin{(-\frac{\pi}{12})})$, b)$-24$, c)$2(\cos{\frac{\pi}{2}} + i \sin{\frac{\pi}{2}})$, d)$\frac{3\sqrt{2}}{2}(\cos{\frac{\pi}{4}} + i \sin{\frac{\pi}{4}})$
\end{description}

\end{document}